\section{Possible Issues}
\textbf{1:} Pen-and-paper method does first-come, first-served.Here urgent patients may be delayed if staff are occupied with low-priority tasks. 
Implementing a priority-based scheduling algorithm in a way that hospitals can easily adapt and use it on the Workflows will aid here.
 \textbf{Example:} While i was working as a doctor, I observed a stable patient with chest pain being deprioritised because their initial ECG was normal. As he was having the chest pain for a week,  CT angiogram was performed as the patient insisted(According to the guideline when the ECG is normal CT angiogram is need).But the CT showed a 90\% blockage requiring emergency  surgery.
\textbf{2:} Lack of optimized data structures and architectural design makes  retrieval and updates delay adoption from the users.
Using  hash maps for quick lookups and priority queues for task scheduling to reduce  latency would help in this case so that  healthcare workers can use the system easily
\textbf{3:} Government hospitals still serve more than half the population but remain under-resourced with no standard digital entry point for EHR data. 
Designing the system to be modular and adaptable making healthcare workers access it directly from their phone will help in this case.
\noindent\textit{These issues form the core focus of this project.}
